\chapter{Details of the sensor tests}
\label{app:sensortests}

The errors in this appendix are given in the definition of the paper (first number) and in the strict
definition (second number, \cref{ch:model}). The strict one shows the differences between settings better. The tests
start from one basic setting: the eight real stations, 5,000 training steps and 3 networks. This setting has an
error of about 6\,\% in the definition of the paper and 44.1\,\% in the strict definition on the first test split,
and the changes in \tref{tab:tried} and \tref{tab:tried2} are compared with these values. A dash means that I did not
measure the value. The sensor in the middle of the room is not part of the basic setting,
because its exact position is not known yet. I tested it as a virtual sensor in the simulation (first rows of
\tref{tab:tried}). The reproduction of the paper has 29.8\,\% in this definition (see the end of this appendix).

Some settings need a short explanation:
\begin{itemize}
  \item \emph{Relative loss}: in training every sample is weighted by the size of its own signal, so that
        scenarios with a small \cotwo{} rise count as much as strong ones. This is how the relative error is
        calculated.
  \item \emph{Tuned}: relative loss and a smaller network (width 64 and 2 blocks instead of width 128 and 3 blocks).
  \item \emph{Anchoring}: first the station values are interpolated to every point of the room. The model then only
        has to predict how much these values change.
  \item \emph{Station dropout}: in half of the training steps some stations are hidden.
  \item \emph{Physics input}: the model also gets the fresh-air flow through the open windows, the \cotwo{} level
        that the room would approach for this flow and this number of people, the time constant of the room, and
        the number of open windows and doors.
\end{itemize}

\section{What I tried}
\begin{table}[h]
  \centering
  \caption[Changes that helped.]{Changes to the sensor model that lowered the error (paper and strict definition).}
  \label{tab:tried}
  \small
  \begin{tabular}{>{\raggedright}p{5.3cm}>{\raggedright}p{3.2cm}>{\raggedright\arraybackslash}p{4.4cm}}
    \toprule
    Change & Paper definition & Strict definition\\
    \midrule
    A grid of 16 virtual sensors & 3.6\,\% & 21.3\,\%\\
    Four virtual sensors & 4.6\,\% & 27.7\,\%\\
    One virtual sensor in the middle & 5.0\,\% & 28.8\,\%\\
    Relative loss and smaller network & 6.0\,\% & 35.1\,\%\\
    Relative loss & 6.0\,\% & 38.8\,\%\\
    More noise, station dropout, stronger weight decay & -- & 39.7, 40.0 and 42.5\,\%\\
    Physics input & 0.0 to 0.1 points better on three splits & about 1 point better on all four splits\\
    \bottomrule
  \end{tabular}
\end{table}

\begin{table}[h]
  \centering
  \caption[Changes without a clear effect.]{Other changes to the sensor model, which made no clear difference.}
  \label{tab:tried2}
  \small
  \begin{tabular}{>{\raggedright}p{5.3cm}>{\raggedright}p{3.2cm}>{\raggedright\arraybackslash}p{4.4cm}}
    \toprule
    Change & Paper definition & Strict definition\\
    \midrule
    Door 2 as the eighth station & 7 stations 6.2\,\%, 8 stations 6.1\,\% & 7 stations 44.2\,\%, 8 stations 44.5\,\%\\
    Anchoring & -- & 44.4\,\%\\
    Longer station history (6, 12 and 20~min) & -- & 40.9, 38.6 and 43.1\,\%\\
    Shorter training & -- & 44.1\,\% instead of 44.2 to 44.5\,\%\\
    \bottomrule
  \end{tabular}
\end{table}

The place of one virtual sensor in the seating area matters little: the middle (28.8\,\%), the door side
(29.8\,\%), the window side (30.2\,\%) and the back (30.4\,\%) are about the same. One place on the blackboard
side gave no gain (43.2\,\%). Training with 13, 26, 39 and 52 scenarios gave 44.6, 47.4, 46.6 and 44.1\,\%, so
more scenarios of the same kind did not help.

\section{Stability on four test splits}
\begin{table}[!htbp]
  \centering
  \caption[Four test splits.]{Error on four test splits (paper and strict definition). The first split is the fixed
  one, the others are random.}
  \label{tab:splits}
  \footnotesize
  \setlength{\tabcolsep}{3pt}
  \begin{tabular}{>{\raggedright}p{3.2cm}rrrrrr}
    \toprule
     & \multicolumn{2}{c}{Default} & \multicolumn{2}{c}{Tuned} & \multicolumn{2}{c}{Tuned with physics}\\
    \cmidrule(lr){2-3}\cmidrule(lr){4-5}\cmidrule(lr){6-7}
    Test split & Paper & Strict & Paper & Strict & Paper & Strict\\
    \midrule
    First split & 6.3\,\% & 44.1\,\% & 6.0\,\% & 35.1\,\% & 6.1\,\% & 34.8\,\%\\
    Random split 1 & 5.6\,\% & 32.3\,\% & 5.8\,\% & 32.5\,\% & 5.7\,\% & 30.7\,\%\\
    Random split 2 & 5.3\,\% & 30.4\,\% & 5.7\,\% & 31.8\,\% & 5.7\,\% & 31.7\,\%\\
    Random split 3 & 6.5\,\% & 32.3\,\% & 6.8\,\% & 31.3\,\% & 6.7\,\% & 30.2\,\%\\
    Mean of random splits & 5.8\,\% & 31.7\,\% & 6.1\,\% & 31.9\,\% & 6.0\,\% & 30.9\,\%\\
    Mean of the four & 5.9\,\% & 34.8\,\% & 6.1\,\% & 32.7\,\% & 6.1\,\% & 31.8\,\%\\
    \bottomrule
  \end{tabular}
\end{table}

On the random splits all settings reach 5.3 to 6.8\,\% in the definition of the paper and about 30 to 32\,\% in
the strict definition. So the first split was a hard one, and the gain of the tuning mostly does not repeat on the
random splits. When I train the tuned model for 20,000 steps with 5 networks, the error in the definition of the
paper stays at 6.0\,\%, and the strict error is 40.0\,\% instead of 35.1\,\%. The model then fits the training
scenarios better the longer it trains. In the definition of the paper the model with the physics input and the tuned
model are within 0.1 points. In the strict definition the model with the physics input is lower than the tuned
model on all four splits, but only by about one point.

For comparison, my reproduction of the paper has 11.5\,\% on the test set in the definition of the paper and
29.8\,\% in the strict definition (\cref{ch:repro}). The full-map model has 0.14\,\% and 0.40\,\%
(\cref{ch:results}). The problems are different, so these numbers are only for orientation.
